\section{Experiments}
\label{sec:experiments}

This section evaluates the proposed training-free loop wrapper across dense and MoE checkpoints, base and instruction-tuned variants, and both standard MHA and MLA backbones, using a fixed recipe with no per-cell hyperparameter tuning (\Cref{sec:setup}), unless otherwise stated.

\paragraph{Main findings.}
Our method yields the largest and most reliable gains on knowledge-heavy multiple-choice tasks, especially MMLU-Pro, GPQA-Main, and ARC-Challenge (\Cref{sec:headline}). Across architectures, most cells are positive or neutral, with failures concentrated in small distilled checkpoints on some knowledge-MC tasks (\Cref{sec:matrix}).
\begin{figure}[t]
\centering
\includegraphics[width=\linewidth]{figures/headline_bars_1.pdf}
\caption{\textbf{Per-benchmark accuracy across three Qwen3 model variants under the training-free loop wrapper.} Each panel shows baseline (striped gray) vs.\ our wrapper (solid blue) on 4 knowledge-MC benchmarks. The per-panel y-axis is cropped to emphasize the baseline-to-loop gap. \textbf{Left:} Qwen3-4B-Instruct (dense MHA) on four mid-range general benchmarks. \textbf{Middle:} Qwen3-4B-Base on four hard MMLU 5-shot subjects, selected from MMLU's $57$ subjects per Appendix~\ref{app:per-subject}). \textbf{Right:} Qwen3-30B-A3B-Instruct (large MoE) on four mid-range benchmarks.}
\label{fig:headline-bars}
\end{figure}

\begin{table}[ht]
\centering
\setlength{\aboverulesep}{0pt}
\setlength{\belowrulesep}{0pt}
\caption{\textbf{Window selection for the looping recipe.} \textbf{(a)}~Window-size sweep on Qwen3-1.7B-Base under forward Euler ($K{=}2$): $n{=}4$ is the sweet spot with a sharp cliff at $n{\geq}6$. \textbf{(b)}~The $n{=}4$ choice generalizes across model scales and families. \textbf{(c)}~Block-mode vs.\ layer-mode iteration on MoE backbones at the canonical loop window $[13,16]$ (block $K{=}3$ / layer $K{=}2$): layer-mode yields +0.5 to +1.7 pp on hard MoE benchmarks. Numbers are $\Delta$pp vs.\ baseline. Bold marks the winning variant per row; ``---'' indicates not run.}
\label{tab:window-combined}
 
\begin{minipage}[t]{0.49\linewidth}
\centering
{\small\textbf{(a) Window-size sweep, Qwen3-1.7B-Base}}\par\vspace{3pt}
\scriptsize
\setlength{\tabcolsep}{4pt}
\renewcommand{\arraystretch}{1.18}
\begin{tabular}{@{}c l r r@{}}
\toprule
\textbf{$n$} & \textbf{Window} & \textbf{16-task} & \textbf{$\Delta$} \\
\midrule
1                  & [14]        & 55.72          & +0.18           \\
2                  & [13,14]     & 55.60          & +0.06           \\
3                  & [13,14,15]  & 55.75          & +0.21           \\
\textbf{4} & [12--15]    & \textbf{56.09} & \textbf{+0.55}  \\
6                  & [11--16]    & 54.72          & -0.82           \\
12                 & [8--19]     & 54.91          & -0.63          \\
28 (all)           & [0--27]     & 27.81          & -27.73          \\
\bottomrule
\end{tabular}
\par\vspace{2pt}
 
\vspace{3pt}
 
{\small\textbf{(b) Looped layer sizes across model scales}}\par\vspace{3pt}
\scriptsize
\setlength{\tabcolsep}{4pt}
\renewcommand{\arraystretch}{1.18}
\begin{tabular}{@{}l r r c@{}}
\toprule
\textbf{Model} & \textbf{$n{=}3$ $\Delta$} & \textbf{$n{=}4$ $\Delta$} & \textbf{Winner} \\
\midrule
Qwen3-0.6B-Base~[\citenum{qwen3}]            & \textbf{+0.57} & +0.22           & $n{=}3$ \\
Qwen3-1.7B-Base~[\citenum{qwen3}]            & +0.21          & \textbf{+0.55}  & $n{=}4$ \\
Qwen3-4B-Base~[\citenum{qwen3}]              & +0.43          & \textbf{+0.85}  & $n{=}4$ \\
Llama-3.2-3B-Instruct~[\citenum{llama3}]     & +0.23          & \textbf{+0.45}  & $n{=}4$ \\
\bottomrule
\end{tabular}
\end{minipage}\hfill
%
\begin{minipage}[t]{0.49\linewidth}
\centering
{\small\textbf{(c) Block-mode vs.\ layer-mode (MoE)}}\par\vspace{3pt}
\scriptsize
\setlength{\tabcolsep}{4pt}
\renewcommand{\arraystretch}{1.18}
\begin{tabular}{@{}l c c@{}}
\toprule
\textbf{Benchmark} & \textbf{block} ($K{=}3$) & \textbf{layer} ($K{=}3$) \\
\midrule
\multicolumn{3}{@{}l}{\textit{Qwen1.5-MoE-A2.7B~[\citenum{qwen15moe}]}}\\
\cmidrule(lr){1-3}
ARC-Challenge & +0.17          & \textbf{+0.85}           \\
CSQA          & +0.16          & \textbf{+0.33}           \\
OBQA          & \textbf{+1.00} & -1.40                    \\
MMLU          & -1.18          & \textbf{\phantom{+}0.00} \\
SciQ          & \textbf{+0.20} & ---                   \\
GPQA-Main     & -1.11          & ---                   \\
\midrule
\multicolumn{3}{@{}l}{\textit{Moonlight-16B-A3B~[\citenum{moonlight}]}}\\
\cmidrule(lr){1-3}
ARC-Challenge & -1.45          & \textbf{+0.51} \\
CSQA          & -0.66          & \textbf{+0.49} \\
OBQA          & -0.20          & \textbf{+1.20} \\
MMLU          & \textbf{+0.74} & ---         \\
SciQ          & \textbf{-0.50} & -0.60          \\
GPQA-Main     & -2.00          & \textbf{+0.90} \\
\bottomrule
\end{tabular}
\end{minipage}
\vspace{-0.1cm}
\end{table}
\paragraph{Robustness.}
The effective loop window follows a stable depth-fraction rule, while MoE models require layer-mode iteration to avoid routing thrash (\Cref{sec:depth}, \Cref{tab:window-combined}). Numerous ablation studies confirm robustness to loop strategy, loop count, window placement, cache handling, and decoding settings (\Cref{sec:ablation-summary}).

\paragraph{Comparisons and transfer.}
Compared with naive inference-time looped transformers, our method avoids collapse and performs best across tested cells (\Cref{sec:method-comparison}, \Cref{fig:method-comparison}). A leakage-free transfer to Qwen3-30B-A3B-Instruct further supports cross-architecture generalization (\Cref{tab:scale-30b}).


\subsection{Experimental setup}
\label{sec:setup}

We evaluate seven Hugging Face checkpoint families spanning dense and MoE backbones, base and instruction-tuned variants, and standard MHA vs.\ Multi-head Latent Attention (MLA): Qwen3 dense $\{0.6,1.7,4\}$B (base \& instruct)~\citep{qwen3}, Qwen3-MoE 30B-A3B-Instruct (\texttt{qwen3\_moe})~\citep{qwen3}, Qwen1.5-MoE-A2.7B-Chat (\texttt{qwen2\_moe}, distilled), Llama-3.2 $\{1,3\}$B-Instruct (distilled from Llama-3.1-8B/70B)~\citep{llama3}, DeepSeek-V2-Lite-Chat (MLA + 64-expert MoE)~\citep{deepseekv2}, and Moonlight-16B-A3B-Instruct (\texttt{deepseek\_v3} family, MLA + MoE)~\citep{moonlight}.

All evaluations use \texttt{lm-eval-harness} 0.4.11 \cite{eval-harness} with \texttt{bfloat16} weights and SDPA attention. For every (model, benchmark) cell, we apply a single \emph{out-of-the-box recipe}: $3$-stage Runge–Kutta at the mid 4 layers, block-mode for dense backbones and layer-mode for MoE backbones, with no per-cell hyperparameter search over position, $K$, or strategy. Per-cell variations in absolute layer indices and in secondary settings (cache, decode) follow mechanically from each architecture's layer count and are logged in Appendix~\ref{app:per-cell} for reproducibility; a fully leakage-free transfer of the recipe \emph{with no per-cell variation of any kind} to a held-out architecture (Qwen3-30B-A3B-Instruct) is reported in Table~\ref{tab:scale-30b}.

\begin{table}[t]
\centering
\caption{\textbf{Knowledge-heavy MC benchmarks with the loop wrapper applied out-of-the-box to each frozen checkpoint.} \emph{No per-cell hyperparameter tuning}. $\Delta$pp is improvement over the no-loop baseline at the same prompt. Per-cell configurations are listed in Appendix~\ref{app:per-cell}, and a fully leakage-free single-recipe generalization check on Qwen3-30B-A3B-Instruct is reported in Table~\ref{tab:scale-30b}.}
\label{tab:headline}
{\small\textbf{(a) Dense backbones}}\label{tab:headline-dense}\par
\small
\setlength{\tabcolsep}{10pt}
\renewcommand{\arraystretch}{1.20}
\setlength{\aboverulesep}{0pt}
\setlength{\belowrulesep}{0pt}
\resizebox{\textwidth}{!}{
\begin{tabular}{l l c >{\columncolor{OursRowBg}}c c}
\toprule
\textbf{Model} & \textbf{Benchmark} & \textbf{Base} & \textbf{Loop (Ours)} & $\boldsymbol{\Delta}$pp \\
\midrule
\multirow{3}{*}{Qwen3-4B-Instruct~[\citenum{qwen3}]}
  & MMLU-Pro 5-shot~[\citenum{wang2024mmlupro}] & 0.5714 & \textbf{0.5979} & \textbf{+2.64} \\
  & GPQA-Main 0-shot~[\citenum{rein2024gpqa}]   & 0.3371 & \textbf{0.3571} & \textbf{+2.01} \\
  & CommonsenseQA 7-shot~[\citenum{talmor2019csqa}] & 0.7887 & 0.7993      & +1.06 \\
\midrule
\multirow{3}{*}{Llama-3.2-3B-Instruct~[\citenum{llama3}]}
  & GPQA-Main 0-shot~[\citenum{rein2024gpqa}]              & 0.2991 & 0.3103 & +1.12 \\
  & MMLU-Pro 5-shot~[\citenum{wang2024mmlupro}]            & 0.3164 & 0.3236 & +0.71 \\
  & MMLU 0-shot~[\citenum{hendrycks2021mmlu}]              & 0.5966 & 0.6039 & +0.72 \\
\midrule
Llama-3.2-1B-Instruct~[\citenum{llama3}]
  & GPQA-Main 0-shot~[\citenum{rein2024gpqa}]              & 0.2790 & 0.2969 & +1.79 \\
\bottomrule
\end{tabular}
}

\vspace{0.6em}
{\small\textbf{(b) MoE backbones}}\label{tab:headline-moe}\par
\small
\setlength{\tabcolsep}{10pt}
\renewcommand{\arraystretch}{1.20}
\resizebox{\textwidth}{!}{
\begin{tabular}{l l c >{\columncolor{OursRowBg}}c c}
\toprule
\textbf{Model} & \textbf{Benchmark} & \textbf{Base} & \textbf{Loop (Ours)} & $\boldsymbol{\Delta}$pp \\
\midrule
\multirow{3}{*}{Qwen1.5-MoE-A2.7B~[\citenum{qwen15moe}]}
  & ARC-Challenge 25-shot~[\citenum{clark2018arc}]    & 0.4829 & \textbf{0.5060} & \textbf{+2.30} \\
  & CommonsenseQA 7-shot~[\citenum{talmor2019csqa}]   & 0.7961 & 0.8133          & +1.72 \\
  & OpenBookQA 0-shot~[\citenum{mihaylov2018obqa}]    & 0.3160 & 0.3320          & +1.60 \\
\midrule
\multirow{2}{*}{Moonlight-16B-A3B~[\citenum{moonlight}]}
  & MMLU 0-shot~[\citenum{hendrycks2021mmlu}]         & 0.6786 & 0.6860 & +0.74 \\
  & OpenBookQA 0-shot~[\citenum{mihaylov2018obqa}]    & 0.3160 & 0.3280 & +1.20 \\
\midrule
DeepSeek-V2-Lite-Chat~[\citenum{deepseekv2}]
  & ARC-Challenge 25-shot~[\citenum{clark2018arc}]    & 0.5794 & 0.5879 & +0.85 \\
\bottomrule
\end{tabular}
}
\end{table}

\subsection{The depth fraction rule}
\label{sec:depth}

Across all eight tested architectures, the optimal window's center sits in a narrow band of fractional depth (Table~\ref{tab:window-combined}, and visualized across nine architectures in Figure~\ref{fig:depth-rule} of Appendix~\ref{app:depth-figure}). For models larger than 1.7B, the optimum lies in the upper half (0.43--0.71, mode $\approx$0.50); for sub-1B models, it shifts earlier (0.25--0.56). We hypothesize that head specialization concentrates in late layers, which the loop must avoid~\citep{belrose2023eliciting,lad2024remarkable,men2024short}; in small or heavily distilled models, the late-layer fraction is larger, so the safe window starts earlier~\citep{takase2021lessons,reid2021subformer}.

\subsection{Layer-mode for MoE}
\label{sec:moe-layer}

On MoE backbones~\citep{deepseekv2,moonlight}, default block-mode iteration causes \emph{routing thrash}, in which each block iteration re-evaluates the gating function on slightly perturbed states, accumulating routing-induced rather than representation-induced changes. This can be fixed by switching to layer-mode iteration, as discussed in Section~\ref{sec:layer-mode} and \cite{csordas2024moeut,bae2024relaxed}. Table~\ref{tab:window-combined} compares the two modes on Qwen1.5-MoE and Moonlight, and layer-mode can be seen to flip most negative cells positive, yielding +0.5 to +1.7 pp on hard MoE benchmarks.

\subsection{Results on knowledge-based MC tasks}
\label{sec:headline}

Our main experimental results concern \emph{knowledge-heavy multiple-choice tasks where the baseline is below ceiling, evaluated with a lenient extractor.} On this class of cells the loop acts as a \emph{frozen-context knowledge refiner}, and the gains are largest and most robust. Figure~\ref{fig:headline-bars} and Table~\ref{tab:headline} report the best loop configuration per model family, separated into dense and MoE backbones.

The largest gains (+2.0 to +2.6 pp) appear on the hardest tasks (MMLU-Pro~\citep{wang2024mmlupro}, GPQA-Main~\citep{rein2024gpqa}, ARC-Challenge~\citep{clark2018arc}) for the strongest MHA backbones. MLA-based MoE models (DeepSeek-V2-Lite, Moonlight) move in the same direction but with 3--4$\times$ smaller magnitudes.



\begin{figure}[t]
    \centering
    \includegraphics[width=\linewidth]{figures/k_ablation_two_panel_with_uniform_legend_right.pdf}
    \caption{\textbf{Effect of loop count $K$ on 16-task macro-average accuracy.} \textbf{(a)} Ours (Algorithm~\ref{alg:decode-loop-layer-detailed}) is stable across $K \in \{1,2,\cdots, 24\}$, while uniform loop~\cite{lys2026inner} (Table~\ref{tab:strategies}) fails to scale after $K{\geq}6$. \textbf{(b)} The naive loop variant degrades monotonically, falling to 37.89\% at $K{=}6$ (-17.71 pp).}
    \label{fig:k_ablation}
\end{figure}
\begin{table}[t]
\centering
\caption{\textbf{Qwen3-30B-A3B-Instruct generalization check at $[22,24]$.} The configuration is fixed in advance from the depth-fraction rule of Section~\ref{sec:depth}, with \emph{no per-cell search or tuning}, i.e.\ the configuration was committed before any cell was scored, so each cell is a fully leakage-free transfer (Appendix~\ref{app:scale30b}). Bold marks the better of two committed solver choices ($K{=}2$ stage Runge–Kutta in Algorithm \ref{alg:rk_beta} vs.\ Heun $K{=}1$, defined in Table~\ref{tab:solver-ablation}).
}
\label{tab:scale-30b}
\setlength{\tabcolsep}{4.2pt}
\renewcommand{\arraystretch}{1.12}
\resizebox{\textwidth}{!}{
\begin{tabular}{@{}lrrrrrrrrr@{}}
\toprule
& \multicolumn{8}{c}{Accuracy / exact-match ($\uparrow$)}
& \multicolumn{1}{c}{PPL ($\downarrow$)} \\
\cmidrule(lr){2-9}\cmidrule(l){10-10}
Method
& ARC-Easy
& HellaSwag
& SciQ
& CSQA
& TruthfulQA
& MMLU-flex
& GPQA-Main
& SuperGPQA
& LAMBADA \\
\midrule
Baseline
& 79.04 & 77.74 & 94.80 & 78.71 & 34.15 & 66.67 & 37.28 & 31.00 & 4.12 \\
\midrule
\rowcolor{OursRowBg}
 $K$-stage RK
& \textbf{79.38} & \textbf{77.93} & \textbf{95.00} & \textbf{79.85}
& \textbf{34.64} & \textbf{67.46} & 37.50 & \textbf{31.70} & \textbf{4.11} \\
$\Delta$ vs. baseline
& +0.34 & +0.19 & +0.20 & +1.14
& +0.49 & +0.79 & +0.22 & +0.70 & -0.01 \\
\midrule
\rowcolor{OursRowBg}
Heun $K{=}1$
& 79.08 & 77.78 & 94.90 & 79.61
& 34.27 & 66.67 & \textbf{37.72} & 31.05 & 4.12 \\
$\Delta$ vs. baseline
& +0.04 & +0.04 & +0.10 & +0.90
& +0.12 & 0.00 & +0.44 & +0.05 & 0.00 \\
\bottomrule
\end{tabular}}
\vspace{-0.5em}
\end{table}

\subsection{Cross-architecture summary}
\label{sec:matrix}

Across the seven model families and a uniform knowledge-MC task suite, we score every cell under the fixed recipe of Section~\ref{sec:setup}, resulting in 27 positive ($\Delta>+0.3$ pp, 60\%) and 12 neutral ($|\Delta|\le0.3$ pp, 27\%) cells. The remaining negatives concentrate in a single regime, namely \emph{sub-3B distilled checkpoints on knowledge MC} (e.g. Llama-3.2-1B at multiple-choice MMLU~\citep{hendrycks2021mmlu}). Even in this regime, GPQA-Main~\citep{rein2024gpqa} \emph{does} flip positive on Llama-3.2-1B (+1.79) at very-early position $[4,7]$, showing the failure boundary is task-dependent rather than absolute. The wrapper produces a positive or neutral signal on \emph{87\% of cells} across dense, MoE, base,
instruct, and distilled checkpoints.

\subsection{Ablations}
\label{sec:ablation-summary}

We systematically conducted numerous ablation studies, with full tables deferred to the appendix. These include the integration strategy and damping schedule (Appendix~\ref{app:strategy-ablation}), loop count $K$ (Figure~\ref{fig:k_ablation}), window width and position (Appendices~\ref{app:failures} and~\ref{app:position-landscape}), block-mode versus layer-mode iteration (Appendix~\ref{app:layer-mode-dense}), KV-cache and decode-time handling (Appendices~\ref{app:cache-robustness} and~\ref{app:wallclock}), and recipe-robustness checks varying window position, strategy, iteration count, and cache choice (Appendices~\ref{app:per-cell} and~\ref{app:search-protocol}). We further report robustness checks, per-subject decompositions, large-scale untuned transfer, and failed-configuration analyses in Appendices~\ref{app:robustness}, \ref{app:per-subject}, \ref{app:scale30b}, and~\ref{app:failures}.

\subsection{Comparison with other looped transformer methods}
\label{sec:method-comparison}

We further isolate the contribution of our method by comparing it against two natural alternatives on the same frozen checkpoints~\citep{dehghani2019universal,yang2025looped,geiping2025scaling,zhu2025ouro,bae2024relaxed,bae2025mor,jeddi2026loopformer,oncescu2026recurrent,tang2026looprpt,saunshi2025reasoning,xu2025expressive,wu2025parallel}. The naive looped transformer of \cite{liu2024looped} collapses on every cell, since the iterates leave the regime the post-loop layers were trained on (Section~\ref{sec:strategy}), and accuracy drops by several
percentage points. In comparison, our method remains within the trained regime, producing the highest accuracy on nearly every cell. Figure~\ref{fig:method-comparison} reports the results on Llama-3.2-3B-Instruct (GPQA-Main~\citep{rein2024gpqa}, MMLU-Pro~\citep{wang2024mmlupro}, MMLU~\citep{hendrycks2021mmlu}) and Moonlight-16B-A3B-Instruct (ARC-C~\citep{clark2018arc}, MMLU, CSQA~\citep{talmor2019csqa}).

\begin{figure}[t]
\centering
\includegraphics[width=\linewidth]{figures/method_comparison_no_parcae_rotated_2.pdf}
\caption{\textbf{Comparison with other looped transformer methods on Llama-3.2-3B-Instruct and Moonlight-16B-A3B-Instruct.} On each backbone we report three knowledge-MC benchmarks under three configurations: \emph{baseline} (no loop, original checkpoint), \emph{naive loop} with $K{=}4$, and \emph{ours} (Algorithm~\ref{alg:rk_beta}).}
\vspace{-1em}
\label{fig:method-comparison}
\end{figure}

